\section{Verification and security boundaries}\label{sec:security}
The certificate has two distinct requirements: an accepted interval must
contain the intended cost change, and the proof transcript must reveal no
more than its declared public statement. The first requirement justifies an
update decision. The second is a leakage-relative property of the proof
system; it does not hide information implied by the public aggregates.
Appendix~\ref{sec:security-experiments} specifies the corresponding experiments.

\subsection{Approved inputs and verifier obligations}
The prover may supply false memberships, moments, compact statistics or
interval endpoints. The auditor therefore obtains the catalog root, deployed
model, quantization, circuit shapes and verification-key hashes independently
of the prover's bundle. Each query authorization specifies the complete
candidate centers and catalog scope. The verifier reconstructs the
center-dependent coefficients and residual radii from these approved inputs.
A prover-supplied model identifier or decision string is insufficient.

Preprocessing must cover every original row exactly once and authenticate
its compact statistics and exact deployed membership. The moment commitment
is fixed before an independent auditor samples a uniform nonzero challenge.
The verifier checks the ordered mask chain and the linked moment proof under
the same roots, model and challenge. Each query must include that linked
aggregate relation and complete, nonduplicated switching-batch coverage.
Signed arithmetic is interpreted only within the checked integer ranges.
For refinement, the verifier also checks the original and derived openings
at the same index, canonical radii, approved centers, distinct corrected
indices and the 16-row budget. Corrections are accepted only with the
complete compact proof backbone.

These requirements assume binding commitments, collision-resistant context
hashing, authentic owner and setup-authority pins, and the knowledge
soundness and zero knowledge of the instantiated argument. Groth16 also
requires an uncompromised trusted setup. The committed observations and
encoder outputs are inputs to the task; their provenance and correctness
are not established by these assumptions.

\subsection{Conditional correctness of the decision}
\begin{proposition}[Conditional soundness of the group decision]
For fixed commitments and an independent uniform nonzero post-commit
challenge in $\mathbb F_p$, suppose all preprocessing and query proofs pass
under the approved keys and complete coverage checks. Except for a
commitment or argument failure, or a moment-fingerprint collision of
probability at most $kD/(p-1)$, the accepted interval in
Eq.~\eqref{eq:group-envelope}, divided by $T^6$, contains the exact quantized
cost change on the approved full catalog. Verified bounded refinement
preserves this containment.
\end{proposition}
\begin{proof}
Binding commitments and complete row coverage fix the original rows,
derived statistics and exact old memberships. For fixed false moments,
a nonzero discrepancy polynomial has degree at most $D$. Its root bound,
with a union bound over groups, gives the fingerprint term. Outside that
event, the linked counts and coordinate sums equal the true group moments;
the checked ranges exclude equality caused only by field reduction.
The aggregate relation therefore gives the exact assigned change in
Eq.~\eqref{eq:group-assigned-change}. Independently reconstructed parameters
and authenticated row statistics give conservative switching bounds.
Complete query coverage and Eq.~\eqref{eq:group-change-identity} then imply
containment of $T^6\Delta$.

A refinement proof authenticates the original row and its compact statistics
at one index, reconstructs canonical endpoints and computes its exact
switching contribution. The monotonicity argument in
Section~\ref{sec:refinement} shows that the corresponding endpoint
corrections remain conservative even if the base proof used a larger valid
row radius. Distinct indices prevent repeated corrections. Applying
Eq.~\eqref{eq:refined-interval} consequently preserves containment.
\end{proof}

A negative upper endpoint certifies strict improvement; a positive lower
endpoint certifies strict worsening. The interval $[0,0]$ certifies equality.
An interval otherwise containing zero leaves the decision unresolved.
Consequently, an inconclusive result is a limitation of decision coverage,
not an accepted incorrect sign. Exact fallback must prove the required
costs on the same catalog root and quantization before resolving that case.

The fingerprint bound applies to one preprocessing session, including its
queries. Repeating preprocessing with new commitments and challenges incurs
the sum of the session bounds. For the BN254 scalar field,
$p\simeq2.1888\times10^{76}$; at $k=5,D=784$, the algebraic collision term is
below $1.8\times10^{-73}$ per session and below $1.8\times10^{-67}$ for a
union bound over $10^6$ such sessions. These values describe the fingerprint
event, rather than the overall cryptographic security level. The benchmark
simulates the independent roles locally; authenticated challenge delivery
and prevention of unaccounted retries are deployment requirements.

\subsection{Declared outputs and aggregate inference}\label{sec:leakage}
Rows, memberships, moments, salts, paths and masks are private witnesses.
The public statement contains metadata, roots, approved centers, masked
preprocessing contributions, the exact assigned-change aggregate $A$,
per-batch switching intervals and the resulting decision. Refinement adds
selected row indices and endpoint corrections. Fallback adds the required
aggregate exact costs. The authorization and termination path, including
proof counts, are also observable. Private diagnostic labels and
certificate flags are excluded from the public export.

Zero knowledge is relative to these outputs. In particular, the assigned
aggregate permits recovery of group statistics under sufficiently revealing
authorized queries. Change only coordinate $d$ of candidate center $g$
from $c^0_{gd}$ to $c^0_{gd}+t$. If both signed probes are authorized and
within range, their outputs satisfy
\begin{equation}\label{eq:probe-leakage}
 A(t)=2t(n_gc^0_{gd}-S_{gd})+n_gt^2.
\end{equation}
For nonzero integer $t$, the auditor can recover
\begin{align}
 n_g&=\frac{A(t)+A(-t)}{2t^2},\\
 S_{gd}&=n_gc^0_{gd}-\frac{A(t)-A(-t)}{4t}.
\end{align}
Repeated coordinate probes recover the sums, and a singleton group's sum
reveals its record. A retained implementation check reproduces this
inference from assigned-change outputs. Witness privacy and hash
commitments do not remove information carried by a valid public answer.

Per-batch intervals and row corrections can reveal additional information
about subcatalogs. The current relation exposes these values so that the
verifier can check and aggregate complete coverage directly. It does not
provide a total-only disclosure interface, differential privacy or adaptive
record non-reconstruction. Externally approved queries restrict who can
request an answer, but the implementation enforces no privacy budget that
would establish those stronger properties. A data owner must authorize the
statistics and the inferences they permit.

\subsection{Implementation evidence and scope}
Fresh public-only replay checks proof acceptance without loading original
rows or private moments. Targeted modifications test rejection of changed
centers, roots, corrections, indices, missing coverage and proof bytes.
These checks support consistency of the implementation with the specified
relations. They do not constitute an independent verifier implementation,
a production challenge service or a composable-security audit. The
mathematical guarantee remains conditional on the stated assumptions,
while the experimental evaluation measures coverage and cost on finite
workloads.
